% =====================================================================
%  Sezione aggiuntiva: stabilità / ZMP
%  Usa i comandi già definiti nel preambolo del report:
%    \Co \Ct \Cth   (sigle dei controllori)
%    \figph{altezza}{etichetta}{didascalia}{cosa deve mostrare}
%    \SI{}{} di siunitx
%  Se nel tuo preambolo manca \figph, sostituisci quel blocco con una
%  figure normale: il resto non dipende da macro nostre.
% =====================================================================

\section{Stabilità: che cosa siamo riusciti a misurare}
\label{sec:zmp}

Il tripode alternato è costruito su un'ipotesi di stabilità \emph{statica}:
in ogni istante tre piedi sono a terra e la proiezione del baricentro deve
cadere dentro il triangolo che formano. Abbiamo voluto verificarla invece di
assumerla, strumentando l'impianto con il calcolo in linea del centro di
pressione e del margine rispetto al poligono d'appoggio.

\subsection{Il criterio, e perché qui coincide con la statica}

Il numero di Froude della marcia nominale è
\begin{equation}
\mathrm{Fr} = \frac{v^2}{g\,h} = \frac{0.12^2}{9.81 \times 0.154} = 0.0095 ,
\label{eq:froude}
\end{equation}
cioè le forze inerziali valgono l'uno percento della gravità. In questo
regime lo ZMP collassa sulla proiezione a terra del baricentro entro la
precisione della misura, e il criterio dinamico degenera in quello statico su
cui l'andatura è costruita. Non ci aspettiamo quindi che questa misura
distingua i controllori: ci aspettiamo che \emph{confermi o smentisca
l'ipotesi} su cui poggia tutto il resto.

Dalle forze di contatto misurate ai sei piedi calcoliamo il centro di
pressione sul piano di riferimento $z = 0$:
\begin{equation}
x_{\mathrm{CoP}} =
\frac{\sum_i \left( x_i F_{z,i} - z_i F_{x,i} \right)}{\sum_i F_{z,i}},
\qquad
y_{\mathrm{CoP}} =
\frac{\sum_i \left( y_i F_{z,i} - z_i F_{y,i} \right)}{\sum_i F_{z,i}},
\label{eq:cop}
\end{equation}
dove la somma corre sulle zampe la cui forza normale supera una soglia di
appoggio. Il \emph{margine} è la distanza con segno dal lato più vicino del
poligono convesso dei piedi attivi, positiva all'interno.

\subsection{Terreno piano: la misura è valida, e lo si può verificare}

Sui task in piano il margine resta positivo per tutta la durata della corsa,
in una banda di \SIrange{0.10}{0.14}{\meter}.

Questo valore non va preso sulla fiducia: è prevedibile dalla sola geometria.
Con le posizioni nominali dei piedi del nostro esapode, i vertici del tripode
distano \SI{0.224}{\meter} e \SI{0.243}{\meter} dall'asse del corpo, e la
distanza dal centro del corpo al lato più vicino del triangolo d'appoggio vale
\SI{0.118}{\meter}. La banda misurata contiene quel valore. Il massimo
osservato, \SI{0.155}{\meter}, resta sotto il raggio inscritto dell'esagono
completo (\SI{0.224}{\meter}), coerentemente con gli istanti di transizione in
cui più di tre zampe sono caricate.

\begin{quote}
La misura in piano riproduce la previsione geometrica. L'ipotesi di stabilità
statica su cui è costruito il tripode è quindi verificata, non assunta.
\end{quote}

\figph{62mm}{fig:zmp-margine}{Margine di stabilità rispetto al poligono
d'appoggio durante la salita della rampa di \SI{8}{\degree} (T4). Fino a
$t \approx \SI{4.5}{\second}$ il robot è in piano e il margine resta nella
banda prevista dalla geometria; dal momento in cui le zampe salgono sulla
rampa compaiono escursioni negative di ampiezza crescente, discusse nel
testo.}{grafico del margine vs tempo su T4, 0--20 s, con evidenziato l'istante in cui il primo piede tocca la rampa}

\subsection{Rampa e ostacoli: le escursioni negative non sono eventi fisici}

Sui terreni non piani il margine calcolato scende sotto zero in istanti brevi
e isolati, con picchi fino a $-\SI{0.11}{\meter}$ verso la fine della salita.
La lettura immediata — il robot perde stabilità per frazioni di secondo — è
sbagliata, e la ragione è di principio.

Il centro di pressione definito dalla~\eqref{eq:cop} è una media dei punti di
contatto pesata con le forze normali $F_{z,i} \ge 0$. Se tutti i contatti
giacciono su uno stesso piano $z = \bar{z}$, la~\eqref{eq:cop} si riduce a
\begin{equation}
x_{\mathrm{CoP}} =
\underbrace{\frac{\sum_i x_i F_{z,i}}{\sum_i F_{z,i}}}_{\text{combinazione convessa degli } x_i}
\;-\;
\bar{z}\,\frac{\sum_i F_{x,i}}{\sum_i F_{z,i}} ,
\label{eq:cop-piano}
\end{equation}
e in equilibrio quasi-statico la risultante tangenziale $\sum_i F_{x,i}$ è
nulla: il secondo termine sparisce e il punto è una combinazione convessa dei
contatti, quindi \textbf{interno al poligono per costruzione}. Un centro di
pressione fuori dal proprio poligono d'appoggio non è un evento raro o
transitorio: è impossibile. Ogni campione negativo del nostro grafico è
perciò un artefatto di misura, non un comportamento del robot.

\paragraph{Origine dell'artefatto.} Il piano di riferimento della~\eqref{eq:cop}
è $z = 0$, cioè il pavimento, ma sulla rampa i contatti non sono più lì. La
rampa parte a $x = \SI{0.66}{\meter}$ e sale di $\tan \SI{8}{\degree} = 0.1405$
al metro, quindi verso la fine della corsa i piedi si trovano circa
\SI{0.2}{\meter} sopra il piano su cui stiamo proiettando. Il termine
$\bar{z}\sum_i F_{x,i}$ della~\eqref{eq:cop-piano} non si annulla più durante
i transitori: ogni risultante tangenziale istantanea viene moltiplicata per
quel braccio e sposta il punto calcolato ben oltre il poligono. In ordine di
grandezza, con $\bar{z} \approx \SI{0.2}{\meter}$ e un peso totale di
\SI{15.6}{\newton}, una risultante tangenziale di circa \SI{8}{\newton} per
pochi millisecondi — valore compatibile con le punte di coppia registrate agli
urti su questo task — produce proprio uno spostamento dell'ordine di
\SI{0.1}{\meter}. In piano lo stesso transitorio è moltiplicato per il solo
raggio del piede, \SI{0.01}{\meter}, ed è invisibile: un fattore venti di
differenza.

\paragraph{Quattro riscontri indipendenti.} La lettura "artefatto" non è una
congettura di comodo: è l'unica compatibile con i dati.

\begin{table}[htbp]
\centering
\caption{Perché le escursioni negative sono attribuite alla misura e non al
robot.}
\label{tab:zmp-evidenze}
\small
\begin{tabular}{@{}p{58mm}p{72mm}@{}}
\toprule
\textbf{Osservazione} & \textbf{Conseguenza} \\
\midrule
Nessuna escursione negativa prima di
$t \approx \SI{4.5}{\second}$, cioè finché il robot è in piano &
L'artefatto compare esattamente quando il piano dei contatti si separa dal
piano di riferimento \\[3pt]
L'ampiezza dei picchi cresce in modo monotono con il tempo &
Cresce con la quota raggiunta, cioè con il braccio $\bar{z}$; la pendenza,
costante, non spiegherebbe una crescita \\[3pt]
Le escursioni sono picchi stretti, non uno scostamento sostenuto &
Coincidono con gli urti all'appoggio, gli unici istanti in cui la risultante
tangenziale non è nulla \\[3pt]
Tutti e tre i controllori completano la salita e il superamento
dell'ostacolo, con assetto limitato &
Un centro di pressione a \SI{0.11}{\meter} fuori da un poligono il cui
margine utile è \SI{0.12}{\meter} avrebbe prodotto un ribaltamento \\
\bottomrule
\end{tabular}
\end{table}

\paragraph{Quello che servirebbe.} La correzione non è sulle forze ma sul
riferimento: il centro di pressione va calcolato \emph{sul piano dei
contatti}, ruotando coordinate e forze in quella terna. Su una superficie
inclinata singola, come la rampa, l'operazione è ben definita e il margine
tornerebbe non negativo per costruzione, diventando una misura utilizzabile.
Sul dosso e sul percorso a ostacoli il problema è più profondo: i piedi
toccano superfici diverse nello stesso istante, non esiste un piano di
contatto unico e il criterio ZMP planare \textbf{non è definito}. Lì servirebbe
il criterio tridimensionale sul cono dei wrench di contatto, che esula da
questo lavoro.

\subsection{Che cosa dichiariamo}

\begin{enumerate}
\item \textbf{In piano la misura è valida e l'esito è positivo.} Il centro di
pressione resta dentro il poligono d'appoggio per tutta la corsa, con un
margine che coincide con la previsione geometrica. L'ipotesi di stabilità
statica dell'andatura a tripode è verificata sperimentalmente.

\item \textbf{Su rampa e ostacoli la misura non è valida, e non la usiamo.}
Le escursioni negative sono un artefatto del piano di riferimento, non un
comportamento del robot; riportarle come perdite di stabilità sarebbe un
errore. Coerentemente, nessuna conclusione di questo lavoro sulla stabilità
poggia su quei campioni.

\item \textbf{La stabilità su terreno irregolare resta però supportata da
misure indipendenti}, che non dipendono dal piano di riferimento: tutti e tre
i controllori completano la salita e il superamento dell'ostacolo, e l'assetto
del corpo resta limitato su ogni task, con \Cth{} che contiene il beccheggio
entro \SI{2}{\degree} dove \Co{} e \Ct{} arrivano a \SIrange{7}{10}{\degree}.
\end{enumerate}

Il valore di questa sezione, dal nostro punto di vista, non sta nel numero che
abbiamo ottenuto ma nella distinzione che ci ha costretto a fare: separare
\emph{il robot è instabile} da \emph{lo strumento non è in grado di dirlo}. È
la stessa disciplina applicata al rumore di fondo nella
Sezione~\ref{sec:method}, e in entrambi i casi la parte utile del risultato è
l'elenco di ciò che non si può dichiarare.
